\chapter{Ateliers de fin d'année et d'oral}
\section{O1 · Transformer un thème en question}
\textbf{S82 · 2 h.} Choisir une situation qui permet un raisonnement du programme, puis rédiger une question précise. Exemples : « Comment choisir un nombre d'essais pour rendre un échec rare ? » ; « Comment un modèle de refroidissement permet-il de prévoir un temps d'attente ? » ; « Pourquoi une moyenne devient-elle plus stable quand l'échantillon grandit ? »
Produire une page contenant la question, les hypothèses, un exemple chiffré, une démonstration accessible et une limite du modèle. Prévoir 20 min de recherche individuelle, 30 min de travail mathématique, 30 min de discussion par groupes, 25 min de rédaction et 15 min de retours.
\section{O2 · Exposer et répondre}
\textbf{S84 · 2 h.} Préparer un exposé de 10 min et des réponses aux questions suivantes : pourquoi ce modèle ? quelle hypothèse est décisive ? que signifie chaque variable ? quel résultat est démontré ? que change une autre valeur du paramètre ?
S'entraîner par groupes de trois : un orateur, un interlocuteur et un observateur. L'observateur relève une étape claire, une étape à justifier et une formulation à améliorer. Changer les rôles. Une vidéo n'est utile que si elle est autorisée et effectivement relue.
\section{O3 · Simulation d'oral}
\textbf{S87 et S93.} Préparer deux questions et réaliser un exposé suivi d'un entretien. Utiliser une grille simple : exactitude mathématique, structure, explication des notations, réponse aux objections, qualité de l'échange. Les modalités finales sont celles des consignes et convocations de la session.
\ifprof\begin{profbox}{Organisation des passages}
Un passage de 20 min occupe un tiers d'une heure : six passages au maximum en 2 h, sans compter les transitions. Avec une classe entière, prévoir plusieurs jurys ou répartir les passages ; ne pas annoncer un oral individuel complet pour tous sur un seul créneau. Faire préciser ce qui est prouvé et ce qui est supposé dans chaque modèle. Une grille interne sert au retour pédagogique ; elle ne remplace pas la grille officielle de la session.
\end{profbox}\fi
